\documentclass[sigconf,nonacm]{acmart}
\usepackage{proyectoACM}


\title{SISTEMA LOCAL DE DETECCIÓN DE DRONES MEDIANTE VISIÓN POR COMPUTADORA}


\author{Camilo Millán, Santiago Castellanos}
\affiliation{%
\institution{Universidad Sergio Arboleda}
\city{Bogotá}
\country{Colombia}
}


\renewcommand{\shortauthors}{Autor et al.}


\areaaccion{Seguridad y vigilancia aérea}
\autoridadconsultada{Andrés Ducuara}


\begin{document}


\maketitle


\section{Información General}


\begin{itemize}
\item \textbf{Institución:} Universidad Sergio Arboleda
\item \textbf{Programa:} Ciencias De La Computación E Inteligencia Artificial
\item \textbf{Fecha:} \today
\end{itemize}


\section{Área de Acción}


\cajainfo{\getareaaccion}


\section{Docente o Autoridad Consultada}


\cajainfo{\getautoridadconsultada}


\section{Planteamiento del Problema}


Colombia enfrenta un incremento sostenido en el uso de drones con fines ilícitos por parte de grupos armados organizados y estructuras criminales. Desde 2020 se han documentado ataques con drones cargados de explosivos contra bases militares y poblaciones civiles en departamentos como Cauca, Norte de Santander y Arauca. A esto se suma el uso de aeronaves no tripuladas para contrabando en centros penitenciarios, vigilancia de rutas de narcotráfico y reconocimiento de operaciones de la fuerza pública.

Los sistemas de detección anti-drones disponibles comercialmente (DroneShield, Dedrone, Rafael Drone Dome) superan los \$100,000 USD, lo que los hace inviables para municipios, bases rurales e instalaciones con presupuesto limitado ---precisamente las zonas más afectadas por estas amenazas.

Se identifica entonces la necesidad de un sistema de detección de drones que sea económicamente accesible, técnicamente viable con hardware convencional y capaz de operar de forma autónoma sin depender de infraestructura en la nube.


\section{Justificación}


El uso de drones con fines ilícitos en Colombia ---ataques con explosivos, contrabando en cárceles y vigilancia de rutas de narcotráfico--- representa una amenaza creciente. Sin embargo, los sistemas de detección disponibles superan los \$100,000 USD, haciéndolos inaccesibles para las zonas más vulnerables del país.

Este proyecto plantea una alternativa de costo cero en infraestructura que, mediante un motor de inferencia propio desarrollado en C++ con ONNX Runtime y OpenCV, permite detectar drones en tiempo real utilizando la webcam de un computador portátil, con alertas automáticas vía Telegram. El modelo de detección, entrenado previamente con la arquitectura YOLOv8n, se exporta a formato ONNX y se consume a través de un pipeline de visión por computadora implementado desde cero, sin dependencia de frameworks de alto nivel como Ultralytics o PyTorch en tiempo de ejecución. Al eliminar la dependencia de servicios en la nube, el sistema opera sin costos recurrentes y con latencia mínima.


\section{Objetivos}


\subsection{Objetivo General}


Diseñar e implementar un sistema local de detección de drones de bajo costo, basado en un motor de inferencia propio desarrollado en C++ que consume un modelo ONNX previamente entrenado con la arquitectura YOLOv8n, capaz de identificar aeronaves no tripuladas en tiempo real mediante webcam y emitir notificaciones automáticas a través de Telegram.


\subsection{Objetivos Específicos}


\begin{enumerate}
\item Exportar el modelo YOLOv8n entrenado (\texttt{best.pt}, mAP50 de 99.25\%) a formato ONNX para su consumo mediante ONNX Runtime en C++, desacoplando la inferencia de cualquier framework de entrenamiento.
\item Implementar desde cero en C++ el pipeline de preprocesamiento de imagen: redimensionamiento con \textit{letterboxing} a 640$\times$640, normalización de píxeles al rango [0,1], conversión de formato HWC a CHW y preparación del tensor de entrada compatible con ONNX Runtime.
\item Implementar desde cero en C++ el pipeline de postprocesamiento: decodificación de las salidas del modelo (coordenadas de bounding boxes y scores de confianza), filtrado por umbral de confianza y supresión de no-máximos (NMS) para eliminar detecciones redundantes.
\item Desarrollar el módulo de captura de video en tiempo real desde la webcam del computador mediante OpenCV en C++, con detección de movimiento por diferencia de frames como prefiltro para optimizar el uso de recursos.
\item Desplegar una infraestructura de backend compuesta por un endpoint REST que reciba los eventos de detección desde el módulo C++, persista los eventos en base de datos y dispare notificaciones en tiempo real a través de un bot de Telegram con imagen anotada y métricas de detección.
\item Validar el sistema mediante pruebas end-to-end midiendo latencia de inferencia en C++, precisión del pipeline custom frente al original de Ultralytics, y estabilidad en operación continua.
\end{enumerate}


\section{Metodología}


El desarrollo del proyecto sigue una metodología incremental organizada en cuatro fases:

\subsection{Fase 0: Exportación del modelo y configuración del entorno (Semana 1)}

Exportación del modelo YOLOv8n entrenado (\texttt{best.pt}, mAP50 de 99.25\%) a formato ONNX mediante el script de exportación de Ultralytics (ejecución única). Se configura el entorno de compilación C++ con CMake, se instalan las dependencias (OpenCV 4.x, ONNX Runtime 1.17+) y se verifica la carga exitosa del modelo \texttt{.onnx} desde C++. Se realizan pruebas de velocidad de inferencia cruda (sin pre/postprocesamiento) para establecer la línea base de rendimiento del equipo.

\subsection{Fase 1: Motor de inferencia en C++ (Semanas 1--2)}

Desarrollo del motor de detección en C++ implementando desde cero:
\begin{itemize}
\item \textbf{Preprocesamiento:} Redimensionamiento con \textit{letterboxing} (preservando aspect ratio con padding gris 114), conversión BGR$\rightarrow$RGB, normalización [0,1], reordenamiento HWC$\rightarrow$CHW y construcción del tensor NCHW float32.
\item \textbf{Inferencia:} Carga de la sesión ONNX Runtime, configuración del proveedor de ejecución (CUDA si disponible, CPU como fallback) y ejecución del modelo sobre el tensor de entrada.
\item \textbf{Postprocesamiento:} Decodificación de la salida del modelo (tensor de forma $1 \times 5 \times 8400$ para una clase), transposición a formato por-detección, filtrado por umbral de confianza, conversión de coordenadas $(cx,cy,w,h)$ a $(x_1,y_1,x_2,y_2)$, re-escalado a las coordenadas originales de la imagen, y aplicación del algoritmo de supresión de no-máximos (NMS) con umbral IoU configurable.
\item \textbf{Captura y visualización:} Captura de video desde la webcam mediante \texttt{cv2::VideoCapture}, detección de movimiento por diferencia de frames (\texttt{cv2::absdiff}) como prefiltro, renderizado de bounding boxes anotados y lógica de cooldown para evitar alertas duplicadas.
\end{itemize}

\subsection{Fase 2: Backend de notificaciones (Semanas 2--3)}

Despliegue de la infraestructura de backend para persistencia y notificaciones. Se implementa un endpoint REST (usando un framework como NestJS o FastAPI) que recibe los eventos de detección desde el módulo C++ vía HTTP POST. El backend persiste cada evento en base de datos (PostgreSQL o SQLite) con la imagen anotada codificada en Base64, el timestamp, el nivel de confianza y metadatos adicionales. Adicionalmente, se integra un bot de Telegram que envía alertas en tiempo real con la imagen de la detección y las métricas asociadas, aplicando políticas de cooldown para evitar saturación.

\subsection{Fase 3: Integración y validación (Semanas 3--4)}

Integración completa del flujo: Webcam $\rightarrow$ Detección de movimiento $\rightarrow$ Preprocesamiento $\rightarrow$ Inferencia ONNX $\rightarrow$ Postprocesamiento (NMS) $\rightarrow$ HTTP POST al endpoint REST $\rightarrow$ Base de datos $\rightarrow$ Telegram. Se realizan pruebas end-to-end del pipeline completo, validación cruzada de las detecciones del motor C++ frente al pipeline original de Ultralytics en Python, medición de latencia en cada etapa, optimización de parámetros (umbrales de confianza, umbral IoU para NMS, intervalo de cooldown, resolución de captura) y prueba de estabilidad en operación continua.


\section{Plan de Trabajo}


\begin{table}[h]
\centering
\begin{tabular}{|c|l|l|}
\hline
\textbf{Semana} & \textbf{Fase} & \textbf{Actividades principales} \\
\hline
1 & Modelo + Entorno & Exportar ONNX, instalar OpenCV + ONNX Runtime, CMake \\
\hline
1--2 & Motor C++ & Pre/postprocesamiento, inferencia, NMS, webcam \\
\hline
2--3 & Backend & Endpoint REST, base de datos, bot Telegram \\
\hline
3--4 & Integración & Pruebas E2E, validación vs. Ultralytics, estabilidad \\
\hline
4 & Ajustes & Optimizar umbrales, documentación final \\
\hline
\end{tabular}
\caption{Cronograma de implementación}
\end{table}


\section{Resultados Esperados}


\begin{itemize}
\item \textbf{Sistema funcional:} Prototipo operativo de detección de drones en C++ que cubra el ciclo completo desde captura de video por webcam hasta notificación al usuario vía Telegram, operando de forma autónoma.
\item \textbf{Motor de inferencia propio:} Pipeline completo de preprocesamiento, inferencia ONNX y postprocesamiento (NMS) implementado desde cero en C++, sin dependencia de Ultralytics o PyTorch en tiempo de ejecución.
\item \textbf{Equivalencia de precisión:} Validación de que el motor C++ produce detecciones equivalentes (mAP50 $\geq$ 95\%) a las del pipeline original de Ultralytics en Python.
\item \textbf{Latencia end-to-end:} Tiempo total desde la captura del frame hasta la recepción de la alerta en Telegram inferior a 2 segundos.
\item \textbf{Costo operativo cero:} Operación del sistema sin costos recurrentes de infraestructura en la nube, utilizando únicamente el hardware existente (computador con webcam).
\item \textbf{Backend de notificaciones:} API REST funcional que persiste eventos de detección en base de datos y dispara notificaciones automáticas vía Telegram con imagen anotada.
\item \textbf{Documentación técnica:} Documento de arquitectura, guía de compilación (CMake) y configuración reproducible del sistema.
\end{itemize}


\end{document}
